\section{Intro}

In this chapter, we look at how we can take a static stabilizer code (the colour code) defined by measurements of weight more than two, and reinterpret it as a dynamic stabilizer code in which all measurements are of weight at most two.
The content herein is based entirely on the paper \textit{Floquetifying the Colour Code}~\cite{Townsend_Teague_2023}.

One reading of this chapter is that we've taken paper \Ccite{Bombin_2024} from PsiQuantum, which Floquetified the surface code (amongst other things!), and done the same thing for the colour code.
The idea at the outset was that, since the colour code can implement the full Clifford group transversally, this process this might lead us to a dynamic stabilizer code in which we could implement the full Clifford group just by adjusting the operation schedule, per \Cref{subsec:induced-logical-cliffords}.
In the end, Davydova et al.\ managed this in \Ccite{Davydova_2024} in a much more satisfactory fashion without using the ZX-calculus.

Nonetheless, what we managed still sounded like a win on the face of it -- high-weight measurements are fiddly to implement without spreading errors, the most common method being to add extra \textit{flag circuitry} to the auxiliary circuits that perform the measurements.
Weight-one and -two measurements, on the other hand, are simple enough; weight-one measurements literally couldn't be simpler, and weight-two measurements have the nice property that implementing them via an auxiliary qubit and controlled Pauli gates is always \textit{fault-tolerant} under circuit-level noise (fault-tolerance defined `Gottesman-style'~\cite{gottesman2024surviving}).
Moreover, some hardware platforms can even implement weight-two Pauli measurements natively.

However, the trade-off is that we must combine more measurement results together to form a detector.
The more measurement results we combine together to form a detector, the more likely the detector is to be violated in any given run of the circuit, the harder the decoder's job is to figure out what faults occurred and fix them, and the more likely we are to end up with a logical error, and hence the wrong output from our logical computation.
Indeed, `big detectors' is essentially the reason the Bacon-Shor code~\cite{Shor1995BaconShor,Bacon_2006} doesn't have a threshold, and one can view the remedies given in \Ccite{gidney2023baconthreshold,Li_2019} as two ways to make the detectors smaller.

That said, one can reasonably object to the premise of the last two paragraphs entirely!
Given a stabilizer code defined by a generating set of stabilizers, some of which have weight more than two, one can point out that almost any circuit in the literature that implements this code will ultimately just use weight-one measurements of auxiliary qubits to perform the high-weight measurements.
Since we now know that circuits are themselves codes, is it really fair to say we've reduced the weights of the required measurements in the code?

I would absolutely agree with this complaint.
These days I mostly view the Floquetification procedure we'll describe here (and its generalisation to all stabilizer codes in Rodatz et al.'s paper~\cite{rodatz2024floquetifying}) as just a rather elaborate syndrome extraction scheme.
The key thing from a QEC perspective is to ensure that throughout the procedure one isn't introducing circuit locations where if an error occurs then it can spread to many more qubits (i.e.\ one should ensure that the procedure is \textit{fault-equivalent}~\cite{rodatz2025faulttoleranceconstruction} -- more on this in \Cref{ch:focked-up}).

I confess that this is something we absolutely did not do during this work!
My excuse is that this was the early days of the Floquet code craze, and it was novel back then to show that such a Floquetification procedure could even be found for a particular code.
As has since been made crystal clear by Rodatz et al.\ in \Ccite{rodatz2024floquetifying}, one can indeed do this for all stabilizer codes, and in a much more sensible `distance-preserving' way.
As such, in hindsight I regard this work as very much being a product of its time; its legacy is more that it was another stepping stone along the path to viewing circuits as codes, and perhaps another stepping stone in the ZX-calculus' march into mainstream QEC, rather than being of much practical use.